% runtime/runtime_operators.tex

\chapter{Runtime Operators}

\label{chap:runtime-operators}

\index{Runtime operators}
\index{Operator!runtime}

Runtime operators define how established engineering objects are
interpreted during execution. They consume upstream concepts from
Foundations and State and map between interpretation layers. They do not
redefine alignment identity, pose2/pose3, metric realization, physical
realization, or representation.

\begin{principle}[Operator-layer principle]
\label{princ:operator-layer-principle}
Within AIM, runtime interpretation is organized through explicit
operators connecting distinct state layers.
\end{principle}

% ============================================================
\section{Operator Overview}
% ============================================================

The following symbols denote families of qualified operations in later
chapters:
\[
\mathcal R,
\qquad
\mathcal D,
\qquad
\mathcal O,
\qquad
\mathcal V.
\]

They do not form one mandatory state ladder. A use may compare a target with a
qualified realization, compose the railway state with \(s(t)\), evaluate a
selected model, express a vehicle-relative result, or apply a purpose-specific
evaluation. Each operation identifies its own inputs, applicability and
provenance.

% ============================================================
\section{Realization Operator}
% ============================================================

\begin{definition}[Realization operator]
\label{def:runtime-realization-operator}
\index{Realization!operator}
The realization operator maps target states to realized states:
\[
\mathcal R :
x_{\mathrm{target}}
\mapsto
x_{\mathrm{real}}.
\]
\end{definition}

For pose3 states,
\[
\mathcal R_{\mathrm{pose3}}:
\mathrm{pose3}_{\mathrm{target}}
\mapsto
\mathrm{pose3}_{\mathrm{real}}.
\]

\(\mathcal R\) is explanatory shorthand for a stated realization relation; it
does not make construction, physical state, maintenance history and
observation one process. A measured state remains an observation with its own
provenance and uncertainty.

% ============================================================
\section{Dynamics Operator}
% ============================================================

\begin{definition}[Dynamics operator]
\label{def:runtime-dynamics-operator-operator-chapter}
\index{Dynamics operator}
For an identified response model \(M\), the dynamics operator maps a qualified
railway state, traversal law and model inputs to a model-specific result:
\[
\mathcal D_M :
(\mathcal S_{\mathrm{rail}}(s(t)),s(t),Q_M)
\mapsto
x_M(t).
\]
\end{definition}

Possible outputs include reduced indicators or modelled responses. Their names
do not establish comfort, safety or admissibility.

A realization-aware use supplies the selected qualified realization as one
input to \(\mathcal D_M\); it does not change the operator's model and
traversal requirements.

% ============================================================
\section{Vehicle Interpretation Operator}
% ============================================================

\begin{definition}[Vehicle interpretation operator]
\label{def:runtime-vehicle-operator}
\index{Vehicle state!interpretation operator}
For an identified vehicle model \(M\), the vehicle interpretation operator
maps the qualified railway state, traversal and vehicle inputs to
vehicle-relative results:
\[
\mathcal V_M :
\left(\mathcal S_{\mathrm{rail}}(s(t)),s(t),Q_M\right)
\mapsto
x_{\mathrm{vehicle},M}(t).
\]
\end{definition}

The qualified input set keeps vehicle type, loading, model fidelity,
provenance and applicability distinct rather than hiding them in one generic
context.

% ============================================================
\section{Operational Evaluation Operator}
% ============================================================

\begin{definition}[Operational operator]
\label{def:runtime-operational-operator}
\index{Operational state!evaluation operator}
The operational operator maps runtime dynamic states to operational
evaluation states:
\[
\mathcal O :
x_{\mathrm{op}}
\mapsto
x_{\mathrm{eval}}.
\]
\end{definition}

\(\mathcal O\) evaluates admissibility, standards compliance, comfort
indicators, and optimization objectives without redefining the upstream
state objects it consumes.

% ============================================================
\section{Operator Composition}

\index{Operator composition}
% ============================================================

Operators may be composed only when output and input types match and
realization, traversal, model, purpose, provenance and applicability remain
attached. There is no universal composition that produces vehicle response,
evaluation or decision from pose3 alone.

% ============================================================
\section{Runtime Abstraction and Implementation}
% ============================================================

\begin{principle}[Runtime abstraction principle]
\label{princ:runtime-abstraction-principle}
Runtime operators define interpretation roles first; detailed mechanics
may be added later as specialized implementations.
\end{principle}

\begin{definition}[Operator role]
\label{def:runtime-operator-role}
An operator role specifies the conceptual transformation performed by an
operator, independent of a particular numerical implementation.
\end{definition}

This distinction keeps AIM structurally clear and implementation
flexible while preserving compatibility with higher-fidelity models.

% ============================================================
\section{Canonical Interpretation}
% ============================================================

\begin{proposition}
Within AIM, runtime interpretation is operator-based.
\end{proposition}

\begin{proposition}
The realization operator \(\mathcal R\) maps target states to realized
states.
\end{proposition}

\begin{proposition}
The dynamics operator \(\mathcal D\) maps railway runtime states to
operational dynamic state trajectories.
\end{proposition}

\begin{proposition}
The vehicle operator \(\mathcal V\) maps railway runtime states to
vehicle-relative interpretation states.
\end{proposition}

\begin{proposition}
The operational operator \(\mathcal O\) maps operational states to
evaluation and admissibility states.
\end{proposition}

% ============================================================
\section{Connection to AIM}
% ============================================================

\begin{summary}
Runtime operators provide the structural connection between established
engineering knowledge and runtime interpretation.

Within this framework, \(\mathcal R\) represents realization,
\(\mathcal D\) runtime dynamic state generation, \(\mathcal V\)
vehicle-relative interpretation, and \(\mathcal O\) operational
evaluation. Their composition yields derived runtime engineering states
without redefining upstream concepts.
\end{summary}
